\documentclass[10pt,a4paper]{amsart}
\usepackage[margin=19mm]{geometry}
\usepackage{amsmath,amssymb,mathtools}
\usepackage[scr=rsfs,cal=euler]{mathalfa}
\usepackage{xfrac}
\usepackage[unicode,hidelinks,bookmarks=false]{hyperref}
\newcommand{\ie}{i.e.\ }
\newcommand{\cf}{cf.\ }
\newcommand{\resp}{respectively\ }
\newcommand{\Subsection}{\subsection}
\providecommand{\nobreakdash}{\nobreak\hbox{-}}
\setlength{\emergencystretch}{2em}
\multlinegap0pt
\usepackage{amssymb}
\usepackage{mathtools}
\usepackage{enumerate}
\usepackage{tikz}
\usepackage{algorithm}\usepackage{algpseudocode}
\usepackage[normalem]{ulem}
\newtheorem{proposition}{Proposition}
\title{\bf Self-dual matroids from canonical curves}
\author{Alheydis Geiger, Sachi Hashimoto, Bernd Sturmfels, Raluca Vlad}
\date{Source: arXiv:2212.05910v2 (CC BY 4.0)}
\begin{document}
\maketitle

\noindent\emph{Source and licence.} \bf Self-dual matroids from canonical curves. Version arXiv:2212.05910v2 (CC BY 4.0), distributed by arXiv under the Creative Commons Attribution 4.0 International licence, http://creativecommons.org/licenses/by/4.0/, as recorded in the arXiv licence metadata for this exact version and shown on the abstract page https://arxiv.org/abs/2212.05910v2. Full licence notice: \texttt{licenses/arxiv-2212.05910-CC-BY-4.0.txt}.\par\smallskip

\noindent\textit{Editorial scope.} Counted unit: the article's proposition prop:containment1 (source0.tex lines 2084--2089). The article's complete original proof of it is delivered with it (source0.tex lines 2091--2102, one \texttt{proof} environment of 12 lines). No mathematics is written, repaired or re-derived: the statement and the proof are the article's own lines, in the article's own notation, and no retained line is substituted or re-flowed.  No further source-local context is needed: every cross-reference of every retained line resolves inside the two delivered spans.  Nothing was omitted from any retained span, so the package carries no omission record.  No retained line contains a citation, so the wrapper carries no bibliography and no bibliography entry.  No retained line contains a figure environment, an image-inclusion command or drawing code, so no image is delivered and the package declares no asset dependency.  Editorial formatting only, in addition: an AMS \texttt{amsart} preamble that restates the article's own declarations and macros used by the retained lines, namely the environments \texttt{proposition} on separate counters, together with the standard packages those lines need.  The retained environments print in this document as Proposition 1 in order of appearance, whereas the article prints the same items with the numbering of its own full document.  The complete native source of the article as distributed in the arXiv e-print for this exact version is bundled unchanged as \texttt{sources/134761/source0.tex}.
\begin{proposition}
  \label{prop:containment1}
  If $C$ and $D$ are circuits in a matroid $M$, with $|C| = l$, $|D| = k$, and $|C \cap D| = l-1$, then any size $k$ subset of $C \cup D$ is dependent in $M$.
  %Suppose that a minimal $\ell$-cycle $C$ of the graph $G$ intersects a minimal $k$-cycle $D$ in $\ell-1$ vertices. 
  %Then any subset of $k$ vertices in $C\cup D$ is dependent in $M_G$.
\end{proposition}

\begin{proof} 
We know that $|C \cup D| = |C| + |D| - |C \cap D| = k+1.$ 
%so $C \setminus D$ is a singleton $\{i\}$. 
So a size $k$ subset $S \subset C \cup D$ is either equal to $D$, which is dependent, or is of the form $(C \cup D) \setminus \{i\}$, for some $i \in D$. In the latter case, if $i \in D\setminus C$, then $S$ contains the circuit $C$, hence is dependent. Otherwise, we must have $i \in C \cap D$; in this case, $S$ contains a circuit by the third circuit axiom, so $S$ is again dependent. 
%The set $\{p_i : i \in C\}$
%spans a projective space $P_C$ of dimension $\leq \ell - 2$,
%while $\{p_j : j \in D\}$ spans a projective
%space of dimension $\leq k-2$.
%Since the two circuits intersect in $\ell-1$ vertices,  $P_C \cap P_D$ contains $\ell-1$ 
%points $p_i$. Because $C$ and $D$ are minimal cycles, these span an $(\ell-2)$-dimensional projective space. This has to coincide with $P_C$. It follows that $P_C\subseteq P_D$. Hence, any $k$ points  from $C\cup D$ are in $P_D$. So, they are
 %dependent in $M_G$.
\end{proof}

\end{document}
